\notepage{81}{Introduction to Green's functions}
{\small\textit{Source reading reminder: Sec.\ 8.2, pp.\ 121--123; the textbook is not identified.}}\par
For a time-independent one-particle problem, let
\[
 H_0(\mathbf r)=-\frac{\hbar^2}{2m}\nabla^2+V_0(\mathbf r),
 \qquad [E-H_0(\mathbf r)]\psi_E^0(\mathbf r)=0.
\]
Define a Green's function by
\[
 [E-H_0(\mathbf r)]G_0(\mathbf r,\mathbf r';E)
 =\delta^{(3)}(\mathbf r-\mathbf r').
\]
We wish to solve
\[
 [E-H_0(\mathbf r)-V(\mathbf r)]\psi_E(\mathbf r)=0.
\]
Insert the delta-function identity into $V\psi_E$:
\begin{align*}
 V(\mathbf r)\psi_E(\mathbf r)
 &=\int\dd^3r'\,\delta^{(3)}(\mathbf r-\mathbf r')
                   V(\mathbf r')\psi_E(\mathbf r')\\
 &=\int\dd^3r'\,[E-H_0(\mathbf r)]G_0(\mathbf r,\mathbf r';E)
                   V(\mathbf r')\psi_E(\mathbf r')\\
 &=[E-H_0(\mathbf r)]\int\dd^3r'\,G_0(\mathbf r,\mathbf r';E)
                   V(\mathbf r')\psi_E(\mathbf r').
\end{align*}
Since $V\psi_E=(E-H_0)\psi_E$, the difference is a homogeneous solution:
\[
 \psi_E(\mathbf r)=f(\mathbf r)+
 \int\dd^3r'\,G_0(\mathbf r,\mathbf r';E)V(\mathbf r')\psi_E(\mathbf r'),
 \qquad(E-H_0)f=0.
\]
For the prescribed incident channel choose $f=c\psi_E^0$. Requiring the
same incident normalization when $V\to0$ fixes $c=1$:
\[
 \boxed{\psi_E=\psi_E^0+G_0V\psi_E.}
\]
Iterating this integral equation gives the formal Born series,
\begin{align*}
 \psi_E&=\psi_E^0+G_0V(\psi_E^0+G_0V\psi_E)\\
 &=\psi_E^0+G_0V\psi_E^0+G_0VG_0V\psi_E^0+\cdots\\
 &=\sum_{n=0}^{\infty}(G_0V)^n\psi_E^0.
\end{align*}
\editorial{Boundary conditions select the homogeneous solution and the
Green's function. A general homogeneous solution need not be a single
multiple of one chosen $\psi_E^0$. The series is formal unless its
convergence or a controlled truncation has been established.}
